\section{The principal term for a nonreal mode}\label{sec:model}

Fix a nonreal root of unity $z$, and put $f_z(n)=z^{\Omega(n)}$.
The proof below treats each such $z$ separately. All constants may
depend on $z$. In particular, no estimate uniform as $z\to1$ is claimed.

\subsection{Classical distribution inputs}

For $\Rea s>1$,
\begin{equation}\label{eq:euler}
 \sum_{n\geq1}\frac{f_z(n)}{n^s}=\zeta(s)^zG_z(s),\qquad
 G_z(s)=\prod_p\frac{(1-p^{-s})^z}{1-zp^{-s}}.
\end{equation}
Branches are continued from this half-plane. The product $G_z$ is
analytic near $1$. The Selberg--Delange theorem of de la Bret\`eche
and Tenenbaum \cite[Theorem 1.2 and equation (1.10)]{BT} gives a
fixed-function expansion to arbitrary fixed logarithmic precision.
Its hypotheses hold since $f_z(p)=z$, $|f_z|=1$, and the higher
prime-power sums converge with any fixed exponent $\sigma>1/2$.

Set $A_z(s)=((s-1)\zeta(s))^zG_z(s)$. The cumulative coefficients
are those of $A_z(s)/s$, divided by $\Gamma(z-j)$. The smooth
density coefficients are those of $A_z(s)$. This follows by
Stieltjes partial summation of a finite cumulative expansion, with
its remainder bounded in absolute value. In detail, if
$A_z(1+t)=\sum a_jt^j$ and $A_z(1+t)/(1+t)=\sum c_jt^j$, then
$a_j=c_j+c_{j-1}$. Differentiating the finite terms gives
\[
 \frac{c_j}{\Gamma(z-j)}+
 \frac{(z-j)c_{j-1}}{\Gamma(z-j+1)}
 =\frac{a_j}{\Gamma(z-j)}.
\]
There is no differentiation of an unspecified error.

We also require the centered progression estimate
\begin{equation}\label{eq:bvz}
 \sum_{q\leq R^\theta}\max_{(a,q)=1}
 \left|\sum_{\substack{n\leq R\\n\equiv a\ (q)}}f_z(n)
 -\frac1{\varphi(q)}\sum_{\substack{n\leq R\\(n,q)=1}}f_z(n)\right|
 \ll_{z,A,\theta}R(\log R)^{-A},\qquad\theta<\tfrac12.
\end{equation}
Here is the verification of the hypotheses needed to apply
Granville--Shao \cite[Theorem 1.1(b)]{GS}.
For a nonprincipal character $\chi$, the twisted Dirichlet series
is $L(s,\chi)^z$ times a product absolutely convergent in
$\Rea s>1/2$. Assumption~\ref{ass:inputs} gives a zero-free
half-plane $\Rea s>7/8$. On each closed smaller half-plane,
Borel--Carath\'eodory applied to discs centered at $2+it$ bounds the
analytic logarithm of $L$ by $O(\log(q(2+|t|)))$. Polynomial bounds
for $L$ follow from character periodicity and partial summation;
the Euler series fixes the logarithm at the centers. The discs avoid
zeros and poles. Extra Euler factors of imprimitive characters are
nonzero there. Thus the twisted series has polynomial growth.

Smooth Mellin inversion can be moved to $\Rea s=15/16$. For
$q\leq(\log R)^D$, it gives arbitrary logarithmic saving after
choosing a sufficient transform decay order. Smooth majorants with
transition width $R(\log R)^{-K}$ remove the cutoff with an error
smaller than any prescribed negative logarithmic power. Character
orthogonality gives the centered Siegel--Walfisz criterion.
For larger $q$, choose $D$ large and use the elementary bounds
$R/q+1$ and $R/\varphi(q)$, treating $q>R$ separately.

The logarithmic derivative coefficients satisfy
$\Lambda_{f_z}(p^j)=z^j\log p$. Its convolution inverse is
$\mu f_z$, whose logarithmic derivative coefficients are the negatives
of these. Both functions belong to the class required by
\cite{GS}. On primes, $f_z=z\ind_{\mathbb P}$, so classical prime
Bombieri--Vinogradov gives the other hypothesis of Theorem 1.1(b).
This proves \eqref{eq:bvz}. Smooth versions follow with the sum over
moduli inside partial summation and a strict exponent margin at
every comparable scale. The principal branch at $s=1$ has not been
removed by this argument.

\subsection{Definition and exact covariance}

Put $\rho=23/50$, $y=\exp(L^\rho)$, and
\[
 n_y=\prod_{p\leq y}p^{v_p(n)},\qquad
 V_y=\prod_{p\leq y}(1-1/p),\qquad
 P_{z,y}(s)=\prod_{p\leq y}(1-zp^{-s}).
\]
For a sufficiently large fixed integer $N$, define
\begin{equation}\label{eq:density}
 b_{z,y,N}(v)=\sum_{j=0}^{N-1}
 \frac{(A_zP_{z,y})^{(j)}(1)}{j!\Gamma(z-j)}
          (\log v)^{z-1-j},\qquad v>1.
\end{equation}
Choose a fixed smooth real $\chi$, between zero and one, equal to
zero on $(-\infty,\eps/4]$ and one on $[\eps/2,\infty)$.
Set $\beta(v)=\chi(\log v/L)b_{z,y,N}(v)$ for $v>1$ and zero
otherwise. Define the principal term on every positive integer by
\begin{equation}\label{eq:model}
 \Bx(n)=\frac{f_z(n_y)}{V_y}\,z^d
       \sum_{\boldsymbol\ell}\frac1\ell
          \beta\left(\frac{n}{n_y\ell}\right).
\end{equation}
This definition uses all primes at most $y$, not an active subset.
The same function therefore occurs for every active set used later.

\begin{lemma}\label{lem:modelsize}
Let $\delta_z=1-\Rea z>0$. On each fixed range $n\leq x^C$,
\begin{equation}\label{eq:modelsize}
 |\Bx(n)|\ll_{N,z,d,\eps,C}L^{-(1-\rho)\delta_z}.
\end{equation}
For every $y$-smooth integer $a$,
$\Bx(an)=f_z(a)\Bx(n)$. The logarithmic exponents in bounds for
any fixed number of smooth derivatives can be chosen independently
of fixed $N$; their constants and thresholds may depend on $N$.
\end{lemma}

\begin{proof}
The identity $(an)_y=an_y$ gives exact covariance, also when $a$ and
$n$ share prime factors. For each fixed $j\geq1$,
\[
 |(\log P_{z,y})^{(j)}(1)|\ll_{j,z}(\log y)^j,\qquad
 |(A_zP_{z,y})^{(j)}(1)|\ll_{j,z}|P_{z,y}(1)|(\log y)^j.
\]
The second bound also holds for $j=0$.
These estimates follow by differentiation of the finite Euler
product and $\sum_{p\leq y}(\log p)^j/p\ll_j(\log y)^j$.
On the cutoff support, successive density terms have a factor
$O_{j,z,\eps}(L^{\rho-1})$. Also
\[
 \frac{|P_{z,y}(1)|}{V_y}
 =\exp\left(\delta_z\sum_{p\leq y}\frac1p+O_z(1)\right)
 \asymp_z L^{\rho\delta_z}.
\]
Multiply this by $L^{\Rea z-1}$ and the bounded harmonic tuple
sum. This proves \eqref{eq:modelsize}. The same calculation applies
to a fixed number of logarithmic derivatives. The finite sequence
of decreasing powers $L^{-(1-\rho)j}$ absorbs their constants for
large $x$ depending on $N$.
\end{proof}

The leading coefficient of \eqref{eq:density} is
$G_z(1)P_{z,y}(1)/\Gamma(z)$. Taking absolute values of all Euler
coefficients before summing would lose the bound just proved.

\subsection{Truncations with square and progression bounds}

Let $D_0=\exp(L^{1/2})$ and $h_0=\lfloor L^{.01}\rfloor$.
For $\eta=1/\log y$, Rankin's inequality gives
\begin{equation}\label{eq:rankin}
 \sum_{P^+(a)\leq y}\frac{a^\eta}{a}\ll L^C,
 \qquad
 \sum_{\substack{P^+(a)\leq y\\a>D_0/y}}\frac1a
 \ll L^C\exp(-\tfrac12L^{.04}).
\end{equation}
If $n_y>D_0$, multiplying its prime factors one at a time finds a
divisor in $(D_0/y,D_0]$. This gives a positive short-divisor
majorant for the discarded integers. Fixed local divisor moments
remain logarithmic, since their local Euler factors are
$1+O(1/p)+O(1/p^2)$. Cauchy's inequality then gives arbitrary
logarithmic square savings as well. On a represented physical list
with product $P$, the Rankin weight adds at most
$P^\eta\leq L^{O_J(1)}$. The saving in \eqref{eq:rankin} absorbs
that cost after $J$ is fixed.

For roughness use the finite Bonferroni expansion
\begin{equation}\label{eq:bonferroni}
 \ind_{P^-(v)>y}
 =\sum_{\substack{e\mid v, P^+(e)\leq y\\
                    \mu^2(e)=1, \omega(e)\leq h_0}}\mu(e)+E(v),
 \qquad |E(v)|\leq\binom{\omega_{\leq y}(v)}{h_0+1}.
\end{equation}
The last expression is a positive divisor sum. Its divisors are
$\leq\exp(O(L^{.47}))$, and their reciprocal mass is at most
$(C\log L)^{h_0+1}/(h_0+1)!$. The retained absolute reciprocal mass
is $L^{O(1)}$.

For a tuple term $u=\ell v$, a divisor condition $a\mid v$ and
compatible congruences $r\mid u+1$, $u+1\equiv b\pmod M$ give
$O_M(x/(\ell ra)+1)$ integers on the support. There are
$O_d(x^{1-\eps})$ tuples. Every fixed product of the short-divisor
supports above has ordinary coefficient mass $x^{o(1)}$. Summing
the rounding errors over $r\leq x^{c_0}$ therefore costs
\begin{equation}\label{eq:rounding}
 x^{1-\eps+c_0+o(1)}=O_A(xL^{-A}).
\end{equation}
For model terms the tuple is harmonic rather than a divisibility
condition, and the rounding cost is only $x^{c_0+o(1)}$.

For square bounds, sum the omitted roughness terms over the possible
smooth parts before applying a norm inequality. Apart from a fixed
logarithmic factor and the bounded harmonic tuple sum, the pointwise
error is at most
\[
 \tau(n_y)2^{\omega_{\leq y}(n)}
       \ind_{\omega_{\leq y}(n)\geq h_0}.
\]
Squaring, using $4^w\ind_{w\geq h_0}\leq2^{-h_0}8^w$, and applying
a fixed divisor moment gives $O(RL^C2^{-h_0})$ on a dyad of scale $R$.
This is the aggregate square bound needed below. Summing separate
norms with a factor $\sum_{a\leq D_0}a^{-1/2}$ would not suffice.

\subsection{Matching the local principal terms}

\begin{lemma}[Rough density]\label{lem:roughdensity}
Fix a positive residual exponent $\eps_1$ and fixed logarithmic
modulus, time, and derivative bounds. For $R\in[x^{\eps_1},x^C]$,
$k\leq L^C$, $(a,k)=1$, $|t|\leq L^B$, and a smooth dyadic test
$\psi$, every sufficiently large fixed $N$ gives
\begin{equation}\label{eq:roughdensity}
 \frac1R\sum_{\substack{v\equiv a\ (k)\\P^-(v)>y}}
 f_z(v)\psi(v/R)(v/R)^{it}
 =\frac1{\varphi(k)}\int\psi(u)u^{it}b_{z,y,N}(Ru)\dd u
       +O_A(L^{-A}).
\end{equation}
For unweighted rough integers the density is $V_y/\varphi(k)$,
with the same accuracy and smoothness bounds.
\end{lemma}

\begin{proof}
Every prime of $k$ is at most $y$ for large $x$.
Use \eqref{eq:bonferroni}. The nonprincipal character sums save
every power of $L$ by the smooth Siegel--Walfisz estimate proved
above, on scales $R/e\geq x^{\eps_1-o(1)}$.
The retained reciprocal mass and the error envelope pay the fixed
modulus and time costs.

For the principal character, apply the fixed-function
Selberg--Delange expansion to each residual scale $R/e$.
Taylor-expand $\log(Ru/e)$ around $\log(Ru)$. Since
$\log e=O(L^{.47})$, each further Taylor order gains a factor
$L^{-.53}$ before fixed costs. The restored moments are
\[
 \sum_{P^+(e)\leq y}\frac{\mu(e)f_z(e)}e(-\log e)^j
       =P_{z,y}^{(j)}(1).
\]
Restoring the omitted degrees in each fixed moment costs at most
\[
 2^{-h_0}\sum_{e\mid\prod_{p\leq y}p}
       \frac{2^{\omega(e)}(\log e)^j}{e}
 \ll_j 2^{-h_0}L^{C_j}.
\]
The derivative product rule and
\[
 \frac1{\Gamma(z-j)}\binom{z-1-j}{r}
       =\frac1{r!\Gamma(z-j-r)}
\]
give \eqref{eq:density}. Choosing sufficiently many fixed terms
gives the desired error; the bounds in Lemma~\ref{lem:modelsize}
permit any fixed $N$ beyond a sufficient threshold.
The principal character is one on every rough integer, so its
contribution is divided by $\varphi(k)$.
For the unweighted case, CRT in \eqref{eq:bonferroni} gives density
$k^{-1}\prod_{p\leq y,p\nmid k}(1-1/p)=V_y/\varphi(k)$.
\end{proof}

\begin{lemma}[Local matching]\label{lem:localmatch}
Let $H_x(n)=H_x(n_y)$, with $|H_x|\leq L^C$. At $Y=x^{1+o(1)}$,
uniformly for logarithmic $k,t$ and the permitted smooth test,
\begin{equation}\label{eq:localmatch}
 \frac1Y\sum_n\psi(n/Y)(n/Y)^{it}H_x(n)
       [f_z(n)T_x(n)-\Bx(n)]\e(hn/k)=O_A(L^{-A}).
\end{equation}
\end{lemma}

\begin{proof}
Discard $n_y>D_0$ using the preceding bounds. In the genuine tuple
term write $n=\ell av$, where $a=n_y$ and $v$ is $y$-rough.
The coefficient is $z^df_z(a)H_x(a)$, and
$R=Y/(\ell a)\geq x^{\eps-o(1)}$. Lemma~\ref{lem:roughdensity}
gives the principal term
\begin{equation}\label{eq:localmain}
 \frac{z^df_z(a)H_x(a)}{\ell a}
 \frac{c_k(ha\ell)}{\varphi(k)}
 \int\psi(u)u^{it}b_{z,y,N}(Yu/(\ell a))\dd u,
\end{equation}
where $c_k$ is the Ramanujan sum over unit residues.
For the model, write $n=av$ and use unweighted rough counting.
Its $V_y$ cancels the $1/V_y$ in \eqref{eq:model}.
Since $\ell$ is a unit modulo $k$, $c_k(ha\ell)=c_k(ha)$, also
when $a$ is a nonunit. The model term is therefore exactly
\eqref{eq:localmain}. The cutoff $\chi$ is one on these residual
scales. Sum the errors with $\sum_a1/a\leq L^C$ and
$\sum_{\boldsymbol\ell}1/\ell=O(1)$, choosing precision after all
test and time exponents. No regularity of $H_x(a)$ is needed.
\end{proof}

\subsection{The discrete high-time estimate}

Take the first comparison core $f_zT_x-\Bx$. Lemma~\ref{lem:localmatch}
handles every fixed logarithmic range of low times. At high times,
extract one small mark. For the genuine tuple contribution, extract
one long prime as well. The identity
$f_z(p_sp_\ell v)=z^2f_z(v)$ leaves the two prime polynomials
unchanged up to phases. Their residual coefficient has unchanged
square norm, so the proof of Lemma~\ref{lem:lambda-correlation}
applies to this part.

For the model contribution, exact covariance gives
$\Bx(p_sn)=z\Bx(n)$. At the remaining scale
$R_0\asymp Y/P_s$, write $n_y=a\leq D_0$ and expand the roughness
of $n/a$. The resulting terms are
\begin{equation}\label{eq:discretemodel}
 \frac{z^df_z(a)H_{\rm rem}(a)\mu(e)}{V_y\ell}
       \ind_{ae\mid n}\,\beta\left(\frac n{a\ell}\right),
\end{equation}
where $e$ is squarefree, $P^+(e)\leq y$, and $\omega(e)\leq h_0$.
After writing $n=aev$, the normalized reciprocal coefficient mass
is $L^{C_*}$. The smooth amplitude is $\beta(ev/\ell)$; its total
variation obeys the same bound, with exponent independent of fixed $N$.
All residual scales $R'=R_0/(ae)$ satisfy $R'\geq x^{1-o(1)}$.

The omitted coefficient has square mass $O_A(R_0L^{-A})$.
Ordinary Dirichlet mean square, applied before restriction to
exceptional times, gives
\begin{equation}\label{eq:modelerror}
 R_0^{-2}\int_{|t|\leq2T_0}\left|\sum_n E_n n^{it}\right|^2\dd t
 \ll L^{-A}(T_0/R_0+\log R_0),\qquad
 T_0/R_0\asymp LP_s/H=o(1).
\end{equation}
The last relation uses both lower bounds on $U_i$ in
\eqref{eq:scales}, not only $Y=x^{1+o(1)}$.

Let $T_{\rm mod}=\exp(L/(\log L)^2)$. For
$L^B\leq|t|\leq T_{\rm mod}$, separate the additive phase into
residue classes modulo $k$. Sum--integral comparison and partial
summation on each progression give a normalized bound
\[
 |N_0(t)|\ll L^{C_*}(|t|^{-1}+x^{-c}).
\]
Indeed the discrete error is $O(L^C(1+|t|)/R')$, while integration
by parts gives the $|t|^{-1}$ term. Consequently
\begin{equation}\label{eq:moderate}
 \int_{L^B\leq|t|\leq T_{\rm mod}}|N_0(t)|^2\dd t
 \ll L^{2C_*-B}+x^{-c'}.
\end{equation}
Logarithmic twists are absorbed by increasing $B$.
For $T_{\rm mod}<|t|\leq2T_0$, the exact ranges in
\eqref{eq:logphase} apply at $R'$ and give
\begin{equation}\label{eq:highmodel}
 |N_0(t)|\ll L^{C_*}\exp(-L/(\log L)^{C_{19}}).
\end{equation}
The residue may be a nonunit; there is no inversion of $ae$.

Outside the exceptional small-prime intervals, the small polynomial
and ordinary mean square on scale $Y/P_s\gg T_0$ suffice.
On the exceptional intervals, integrate \eqref{eq:moderate}
directly and use \eqref{eq:highmodel} over at most
$\exp(O(L^{.9}))$ unit intervals. The latter loss is absorbed by
its exponential saving. This proves \eqref{eq:energy} for the
model part. Choose $B$ after the fixed kernel, group, and derivative
exponents, then choose $N$ after $B$ and the low-time accuracy.
The derivative exponents in Lemma~\ref{lem:modelsize} make this
order possible.

\begin{proposition}[Marked model transfer]\label{prop:modeltransfer}
For the data of Section~\ref{sec:framework}, every fixed $M,r$ and
every required accuracy $A$, all sufficiently large fixed $N$ satisfy
\begin{equation}\label{eq:modeltransfer}
 \sum_{w\equiv r\ (M)}\frac{\Psi(w/x)}w\W(w)
 [f_z(w)T_x(w)-\Bx(w)]G(w+1)W_1(w)W_1(w+1)
 \ll_A L^{-A}.
\end{equation}
One $N$ can be used for every active set in a fixed extraction.
\end{proposition}

\begin{proof}
The first physical core is eventually bounded, so
\eqref{eq:listnorm} gives a norm exponent independent of $J,N$.
It remains list-independent. Use the original ideal family.
On stripping shared labels, both terms have the exact covariance
\[
 [f_zT_x-\Bx](D_Rw)
      =z^{Js_P-c_\nu}[f_zT_x-\Bx](w).
\]
Divide out $z^{Js_P}$, restore the dyads and Mellin integral, and
apply \eqref{eq:orbit}. The residual contains the raw difference
and all comparison differences with factors $z^{-c_\nu}$.
The local and high-time estimates just proved bound every comparison,
including fixed residue restrictions through Fourier shifts.
Choose their accuracy after the total logarithmic pattern cost and
subtract them. The same $\Bx$ and a sufficiently large single $N$
work for every active set, because their exponents have a common
fixed bound in each extraction.
\end{proof}
